\section{Mediator}
\label{sec:pat-mediator}

\problemstatement{Define an object that encapsulates how a set of objects
interact, so they refer to the mediator instead of to each other -- turning
a tangle of many-to-many references into a hub-and-spoke.}

\begin{patternbox}[When You Reach for It]
The trigger is \emph{``every component talks to every other component''} --
an $n^2$ web of direct references (a dialog where each widget reacts to
others, a chat room, air-traffic control). When coupling between peers
grows quadratic, route it through a mediator.
\end{patternbox}

\subsection*{Understanding the pattern}

Instead of colleagues holding references to one another, each holds a
reference to a single \emph{mediator}. When a colleague changes, it
notifies the mediator, which decides who else needs to know and calls
them. The colleagues become simple and reusable; all the interaction logic
concentrates in one place -- the mediator -- which is the only object that
knows the participants.

\begin{ideabox}[Centralise the Many-to-Many]
Replace peer-to-peer links with a hub. Colleagues know only the mediator;
the mediator knows the colleagues. This converts $O(n^2)$ coupling into
$O(n)$ and makes interaction rules explicit and changeable in one spot.
\end{ideabox}

\subsection*{C++ implementation}

\begin{lstlisting}
#include <bits/stdc++.h>
using namespace std;

class User;
struct ChatRoom {                               // mediator
    virtual void send(const string& msg, User* from) = 0;
    virtual ~ChatRoom() = default;
};

class User {
    string name;
    ChatRoom* room;                             // knows only the mediator
public:
    User(string n, ChatRoom* r) : name(move(n)), room(r) {}
    const string& getName() const { return name; }
    void send(const string& msg);
    void receive(const string& msg, const string& from) {
        cout << name << " got from " << from << ": " << msg << "\n";
    }
};

class Room : public ChatRoom {
    vector<User*> users;
public:
    void join(User* u) { users.push_back(u); }
    void send(const string& msg, User* from) override {
        for (User* u : users)
            if (u != from) u->receive(msg, from->getName());   // hub routes
    }
};

void User::send(const string& msg) { room->send(msg, this); }

int main() {
    Room room;
    User a("Alice", &room), b("Bob", &room), c("Carol", &room);
    room.join(&a); room.join(&b); room.join(&c);
    a.send("hi all");                           // no direct peer links
    return 0;
}
\end{lstlisting}

\begin{complexitybox}[Trade-offs]
\textbf{Pros.} Removes tight peer-to-peer coupling; interaction logic
lives in one place; colleagues become reusable. \textbf{Cons.} The
mediator can swell into a complex ``god object'' as rules accumulate --
you have centralised complexity, not removed it.
\end{complexitybox}

\begin{takeawaybox}
Mediator replaces a many-to-many mesh of collaborators with a hub that
owns the interaction rules, cutting coupling from $O(n^2)$ to $O(n)$.
Watch the mediator itself for bloat. Classic examples: chat rooms, UI
dialog coordination, air-traffic control.
\end{takeawaybox}
